\documentclass[a4paper]{article}
% Español (México) con XeLaTeX: fontspec para fuentes Unicode, polyglossia
% para la división silábica y los nombres del documento en la variante mexicana.
\usepackage{fontspec}
\usepackage{polyglossia}
\setdefaultlanguage[variant=mexican]{spanish}
% Los nombres chinos van como «pinyin (original)»: xeCJK da a los caracteres
% Han su propia fuente; sin ella XeLaTeX los omite con un aviso.
% FandolSong no cubre algunos caracteres de nombres (U+73FA); WenQuanYi Zen Hei sí.
\usepackage[AutoFallBack=true]{xeCJK}
\setCJKmainfont{FandolSong-Regular.otf}
\setCJKfallbackfamilyfont{\CJKrmdefault}{WenQuanYi Zen Hei}
% polyglossia no traduce los nombres de listings.
\AtBeginDocument{\providecommand{\lstlistingname}{}\renewcommand{\lstlistingname}{Listado}%
  \providecommand{\lstlistlistingname}{}\renewcommand{\lstlistlistingname}{Índice de listados}}
\usepackage{amsmath, amssymb}
\usepackage{graphicx}
\usepackage[margin=2.25cm]{geometry}
\usepackage[most]{tcolorbox}
\usepackage{booktabs}
\usepackage{tabularx}
\usepackage{array}
\usepackage{float}
\usepackage{xcolor}
\usepackage{hyperref}
\usepackage{fancyhdr}
\setCJKmainfont[AutoFakeBold=2.4]{AR PL UMing CN}
\setCJKsansfont[AutoFakeBold=2.4]{Droid Sans Fallback}
\setCJKmonofont[AutoFakeBold=2.4]{Droid Sans Fallback}
\hypersetup{colorlinks=true, linkcolor=blue!55!black, urlcolor=blue!60!black}
\graphicspath{{./}{figures/}}
\setlength{\parskip}{0.55em}
\setlength{\parindent}{2em}
\setlength{\emergencystretch}{3em}
\renewcommand{\arraystretch}{1.18}
\newtcolorbox{knowledgebox}[1]{enhanced, colback=blue!5!white, colframe=blue!70!black, colbacktitle=blue!70!black, coltitle=white, fonttitle=\bfseries, title={#1}, sharp corners, breakable}
\newtcolorbox{importantbox}[1]{enhanced, colback=yellow!10!white, colframe=yellow!75!black, colbacktitle=yellow!75!black, coltitle=black, fonttitle=\bfseries, title={#1}, sharp corners, breakable}
\newtcolorbox{warningbox}[1]{enhanced, colback=red!5!white, colframe=red!70!black, colbacktitle=red!70!black, coltitle=white, fonttitle=\bfseries, title={#1}, sharp corners, breakable}
\pagestyle{fancy}
\fancyhf{}
\fancyfoot[C]{\thepage}
\fancyfoot[R]{\footnotesize\color{gray}\href{https://github.com/hqhq1025/ai-course-notes}{AI Course Notes}}
\renewcommand{\headrulewidth}{0pt}
\renewcommand{\footrulewidth}{0pt}
\newcommand{\videourl}{https://www.youtube.com/watch?v=qtugoE1xQZk}
\newcommand{\repourl}{https://github.com/hqhq1025/ai-course-notes}

\begin{document}
\begin{titlepage}
\centering
{\Large Notas didácticas de una entrevista en profundidad\par}
\vspace{1.0cm}
{\huge\bfseries Diez años de evolución de la conducción autónoma: de los mapas de alta precisión y BEV a los modelos de extremo a extremo y los modelos del mundo}\par
\vspace{0.5cm}
{\Large Zhang Xiaojun Podcast: conversación con Lang Xianpeng (郎咸朋), de Li Auto (理想汽车)\par}
\vspace{0.35cm}
{\large Elaborado a partir del video público, la estructura de capítulos, la transcripción, la descripción del video y diagramas conceptuales propios\par}
\vspace{0.3cm}
{\large 2025-03-16, entrevista grabada en 2024-12\par}
\vspace{0.85cm}
\includegraphics[width=0.88\textwidth,height=0.42\textheight,keepaspectratio]{cover.jpg}\par
\vfill
\begin{tcolorbox}[width=0.92\textwidth, colback=black!2!white, colframe=black!55, sharp corners]
\textbf{Título del video}: 96. Conversación con Lang Xianpeng: diez años de evolución de la conducción autónoma, detalles técnicos clave y Tesla\par
\textbf{Canal del video}: Zhang Xiaojun Podcast\par
\textbf{Duración del video}: 02:00:20\par
\textbf{Enlace del video}: \href{\videourl}{\nolinkurl{\videourl}}\par
\textbf{Nota sobre los materiales}: YouTube no ofrece subtítulos descargables; el texto principal se elaboró a partir de una transcripción por bloques con faster-whisper large-v3 (CPU, int8), la descripción del video, la estructura de capítulos y figuras didácticas propias. Las tomas fijas de la entrevista no se reutilizan en el texto principal.\par
\textbf{Página del proyecto}: \href{\repourl}{\nolinkurl{\repourl}}\par
\end{tcolorbox}
\end{titlepage}

\begingroup
\small
\setlength{\parskip}{0pt}
\renewcommand{\baselinestretch}{0.92}\selectfont
\tableofcontents
\endgroup
\newpage
\section{Introducción: por qué este episodio es una clase sobre las rutas tecnológicas de la conducción autónoma}

Esta sección sitúa el episodio. Lang Xianpeng (郎咸朋) trabajó de 2013 a 2018 en Baidu (百度) en conducción autónoma y mapas de alta precisión; en 2018 se incorporó a Li Auto (理想汽车), y durante la última década ha estado en la primera línea de la conducción autónoma en China. El valor de esta entrevista no está en los chismes sobre una empresa, sino en que, desde la perspectiva de alguien que lo vivió, recorre la ruta tecnológica de la conducción autónoma desde «mapas de alta precisión + LiDAR + reglas» hasta «BEV+Transformer», «de extremo a extremo» y «VLM + modelo del mundo + RL».

\begin{importantbox}{Tesis central del episodio}
El hilo conductor de diez años de conducción autónoma es el paso de desarrollar el vehículo como si fuera un «tranvía» a tratar la conducción como una capacidad que se puede aprender. La ruta temprana dependía de mapas de alta precisión, LiDAR costosos e ingeniería de reglas; la contribución clave de Tesla fue resolver el problema elevando la dimensión con visión, BEV, chips y datos de flota; el enfoque de extremo a extremo comprime además las reglas manuales en el modelo y los datos, y empieza a pasar de funciones de software a un sistema de capacidades.
\end{importantbox}

\begin{knowledgebox}{Nota sobre la estrategia visual}
Este video es una entrevista con encuadre fijo, sin proyección de pantalla ni pizarrón. El texto usa la portada solo para identificar la fuente; todas las imágenes del texto son diagramas conceptuales de elaboración propia, que explican la evolución de la pila tecnológica de la conducción autónoma, la elección de sensores, el enfoque de extremo a extremo, los límites de responsabilidad de L3/L4, el ciclo cerrado de datos de Li Auto y la presión organizacional.
\end{knowledgebox}

\begin{figure}[H]
\centering
\includegraphics[width=0.96\textwidth]{av-ten-year-map.png}
\caption{Mapa de diez años de la conducción autónoma: de los mapas de alta precisión y el LiDAR a BEV/Transformer, el enfoque de extremo a extremo y los modelos del mundo. Diagrama conceptual de elaboración propia, basado en la conversación de 00:01:32--00:41:14.}
\end{figure}

\begin{knowledgebox}{Lectura de la figura: de la certeza de ingeniería a la capacidad del modelo}
La figura va de HD Map y LiDAR Stack a BEV, End-to-End y World Model. Al leerla hay que fijarse en el cambio de paradigma: al principio se usaban mapas y sensores para simplificar el mundo; después, datos y modelos para aprender directamente un mundo complejo; más adelante, el sistema necesita predecir el futuro y optimizarse en ciclo cerrado.
\end{knowledgebox}

\subsection{Ruta de lectura}

Esta sección propone cómo leer el episodio. Todo el video puede organizarse en torno a cinco preguntas: primera, ¿por qué la ruta de mapas de alta precisión y LiDAR pesado se parece a un «tranvía»? Segunda, ¿por qué Tesla insiste en la visión y en chips propios? Tercera, ¿qué problema de modelado espacial resolvieron BEV y Transformer? Cuarta, ¿por qué el enfoque de extremo a extremo implica que «antes todos hacíamos mal la conducción autónoma»? Quinta, ¿cómo determinan en conjunto la siguiente etapa L3, L4, los modelos del mundo, RL y el ciclo cerrado de datos y organización de Li Auto?

\noindent\begin{tabularx}{\textwidth}{p{0.19\textwidth}p{0.34\textwidth}X}
\toprule
Pregunta de lectura & Material de la entrevista & Juicio que hay que formarse \\
\midrule
En qué se equivocó la ruta temprana & Mapas de alta precisión, LiDAR, reglas, enumeración exhaustiva de ODD & Trataba las vías abiertas con una mentalidad de rieles y difícilmente alcanza scale. \\
Por qué importa Tesla & Visión pura, BEV, chips, solución por elevación de dimensión & Llevó la conducción autónoma de la acumulación de sensores a una ruta de datos y modelos. \\
Qué cambia el enfoque de extremo a extremo & Los escenarios no pueden enumerarse, las reglas interfieren entre sí, el modelo aprende capacidades & Pasar del desarrollo de funciones al entrenamiento de capacidades. \\
Cómo entender L3/L4 & Responsabilidad del sistema, toma de control, zonas seguras, ODD & L3 no es una extensión de L2; se parece más a un precursor de L4. \\
Cuál es la experiencia de Li Auto & Datos, cómputo, presión organizacional, proyecto Acrópolis (卫城) & La conducción autónoma en producción en serie es un acoplamiento de tecnología, datos, cadena de suministro y organización. \\
\bottomrule
\end{tabularx}

\begin{warningbox}{Límite temporal}
El video se publicó en marzo de 2025, pero la entrevista se realizó en diciembre de 2024. Cuando en la entrevista se dice «este año» suele referirse a 2024, y «el año pasado», a 2023. El texto interpreta según el contexto original de la entrevista y no confunde los tiempos relativos con la fecha de publicación.
\end{warningbox}

\subsection{Resumen de la sección}

EP96 es uno de los episodios más técnicos de la serie de IA e internet de Zhang Xiaojun (张小珺). Puede leerse junto con EP120 sobre Physical AI de XPeng (小鹏), EP132 sobre Galaxea (星海图) y la versión de VLA con proyección de pantalla: la conducción autónoma no es solo una historia de la industria automotriz, sino la avanzada de la IA del mundo físico, los ciclos cerrados de datos, los modelos de extremo a extremo y los modelos del mundo.
\section{Primera etapa: mapas de alta precisión, LiDAR y la mentalidad del «tranvía»}

El capítulo anterior trazó la hoja de ruta; este capítulo regresa primero al punto de partida, en 2014--2015. Lang Xianpeng (郎咸朋) explica que en aquel momento muchos equipos desarrollaban la conducción autónoma como si fuera un «tranvía»: primero se usaba un mapa de alta precisión para tender un riel virtual sobre la vía real, después se recurría a numerosos LiDAR para detectar los objetos dinámicos del entorno y, por último, se empleaban if-else y sistemas de reglas para controlar el vehículo a lo largo del riel. Hoy parece risible, pero entonces era una decisión de ingeniería muy natural, porque reducía un problema de mundo abierto a un problema relativamente determinista.

\begin{figure}[H]
\centering
\includegraphics[width=0.94\textwidth]{hdmap-lidar-vs-vision.png}
\caption{Mapa de alta precisión + LiDAR frente a la ruta de visión: una depende de la certidumbre externa; la otra, de los datos y de la generalización del modelo. Diagrama conceptual propio, elaborado a partir de la conversación de 00:01:32--00:04:30 y 00:25:06--00:28:46.}
\end{figure}

\begin{knowledgebox}{Lectura de la figura: las dos rutas tratan la incertidumbre de forma distinta}
A la izquierda, la ruta de mapa de alta precisión + LiDAR delega la incertidumbre en el mapa, la medición de distancias y las reglas; a la derecha, la ruta de visión + modelo deja la incertidumbre a los datos y al aprendizaje del modelo. La primera es controlable a corto plazo, pero su mantenimiento es pesado a largo plazo; la segunda es difícil de entrenar a corto plazo, pero a largo plazo ofrece más posibilidades de escalar.
\end{knowledgebox}

\subsection{Por qué el mapa de alta precisión se parece a un riel virtual}

Esta sección explica la metáfora del «tranvía». El mapa de alta precisión registra de antemano las líneas de carril, los límites de la vía, la posición de los semáforos, la topología de las intersecciones y otros datos, lo que equivale a darle al vehículo un riel predecible. Mientras circula en tiempo real, el vehículo no necesita comprender el mundo por completo a partir de sus sensores: alinea su posición actual con el mapa y luego avanza según las reglas. Este método funciona con relativa facilidad en vías cerradas y que cambian poco, como las autopistas, pero en las calles urbanas comunes su costo de mantenimiento es altísimo.

\begin{warningbox}{El problema de fondo del mapa de alta precisión es la actualización y la cobertura}
Las autopistas tienen un kilometraje limitado y cambian poco, por lo que se prestan al uso de mapas de alta precisión; las calles urbanas comunes son muchísimas, tienen obras frecuentes y muchos cambios temporales, así que es difícil actualizar el mapa cada día. Si la conducción autónoma depende de un mapa que siempre sea correcto, los cambios del mundo real terminarán por quebrarla.
\end{warningbox}

\subsection{Por qué acumular LiDAR no permite la producción en serie}

Esta sección examina el costo del hardware en los primeros años. En la entrevista se menciona que, al principio, un LiDAR de 64 líneas podía costar entre 500,000 y 600,000 yuanes, y que los vehículos de ingeniería de Baidu o Cruise llevaban siete u ocho LiDAR cada uno, de modo que el costo de los sensores superaba con mucho el del propio automóvil. Un sistema así sirve para demostraciones de investigación y pruebas de Robotaxi a pequeña escala, pero no para cientos de miles o millones de vehículos producidos en serie.

\begin{figure}[H]
\centering
\includegraphics[width=0.86\textwidth]{sensor-cost-stack.png}
\caption{Pila de costos de la conducción autónoma: las primeras soluciones quedaron lastradas por el costo de los sensores, la plataforma de cómputo y los vehículos de ingeniería. Diagrama conceptual propio, elaborado a partir de la conversación de 00:12:07--00:13:09.}
\end{figure}

\begin{knowledgebox}{Lectura de la figura: el vehículo de investigación y el de producción en serie no tienen las mismas restricciones}
La figura va de los múltiples LiDAR, el conjunto de sensores, la plataforma de cómputo y el vehículo de ingeniería hasta las restricciones de la producción en serie. Un vehículo de investigación puede cambiar sensores costosos por certidumbre; un vehículo de producción en serie debe tomar en cuenta el costo, la apariencia, el consumo de energía, la cadena de suministro, el mantenimiento y el precio para el usuario.
\end{knowledgebox}

\subsection{Términos clave: ODD, mapa de alta precisión y LiDAR}

Esta sección completa los conceptos básicos. En las discusiones sobre conducción autónoma, muchos términos parecen comunes, pero si no se explican, más adelante será difícil seguir L2/L3/L4 y la conducción de extremo a extremo.

\begin{knowledgebox}{Términos clave: las primeras rutas de la conducción autónoma}
\noindent\begin{tabularx}{\textwidth}{p{0.20\textwidth}p{0.36\textwidth}X}
\toprule
Término & Significado & Papel en este episodio \\
\midrule
ODD & Operational Design Domain, la zona y las condiciones en las que el sistema está diseñado para operar & La conducción autónoma temprana intentaba enumerar exhaustivamente los límites de los escenarios mediante el ODD. \\
Mapa de alta precisión & Información de vías, carriles, infraestructura de tránsito y topología con precisión centimétrica & Proporciona un riel virtual, pero su cobertura y actualización son difíciles. \\
LiDAR / radar láser & Sensor que emite activamente un láser y mide distancias & Da la distancia geométrica de forma directa, pero su costo y su densidad de información son limitantes. \\
Sistema de reglas & If-else y lógica de escenarios escritos a mano por ingenieros & Es interpretable, pero no puede enumerar exhaustivamente los escenarios de cola larga. \\
\bottomrule
\end{tabularx}
\end{knowledgebox}

\subsection{Resumen de la sección}

La conducción autónoma temprana eligió el mapa de alta precisión y el LiDAR porque así podía convertir temporalmente la vía abierta en un sistema determinista. Sin embargo, esta ruta chocó con un límite en la cobertura urbana, la actualización de mapas, el costo y los escenarios de cola larga. La ruta de Tesla, en el capítulo siguiente, es precisamente una sustitución de nivel superior de esta ruta de ingeniería determinista.
\section{La ruta de Tesla: visión, BEV, Transformer y chips propios}

La sección anterior expuso los problemas de las primeras rutas; esta sección se ocupa de Tesla. Lang Xianpeng (郎咸朋) considera que hacia 2018 la conducción autónoma dio un giro decisivo: Tesla se mantuvo firme en no usar mapas de alta precisión ni LiDAR y se orientó hacia la visión pura, BEV, Transformer y chips de desarrollo propio. Lo esencial no es ninguno de esos términos por separado, sino llevar la información de varias cámaras a una misma representación espacial para que el sistema comprenda de forma unificada el mundo que rodea al vehículo.

\begin{figure}[H]
\centering
\includegraphics[width=0.96\textwidth]{tesla-updimension.png}
\caption{La solución de Tesla mediante elevación de dimensión: resolver problemas del sistema con datos de la flota, visión, chips y entrenamiento a gran escala. Diagrama conceptual propio, elaborado a partir de la conversación de 00:04:30--00:25:06.}
\end{figure}

\begin{knowledgebox}{Lectura de la figura: elevar la dimensión no es acumular más parches}
En la figura, Fleet Data, Vision, Chip, Training y FSD forman una sola línea. La «elevación de dimensión» de Tesla consiste en dejar de corregir cada error de fusión tardía y, en cambio, modificar la forma de representación y los límites del sistema: resolver el problema con un espacio unificado, un modelo unificado, datos de la flota y capacidad de cómputo propia.
\end{knowledgebox}

\subsection{Por qué la visión}

Esta subsección explica la ruta de visión. Elon Musk suele justificar la visión pura con el argumento de que «las personas manejan solo con los ojos», lo cual funciona para la divulgación; en el plano técnico, lo más importante es que las imágenes tienen una alta densidad de información: contienen color, textura, semántica, bordes y geometría implícita. Mediante la paralaje entre varias cámaras, el movimiento entre fotogramas consecutivos y el aprendizaje del modelo, el sistema puede inferir a partir de las imágenes el espacio tridimensional, las velocidades y las tendencias.

\begin{figure}[H]
\centering
\includegraphics[width=0.92\textwidth]{camera-lidar-tradeoff.png}
\caption{La disyuntiva entre Camera y LiDAR: LiDAR proporciona la distancia de forma directa; Camera aporta semántica y escala a bajo costo. Diagrama conceptual propio, elaborado a partir de la conversación de 00:25:06--00:28:46.}
\end{figure}

\begin{knowledgebox}{Lectura de la figura: medir la distancia directamente no equivale a tener más información}
La ventaja de LiDAR es la precisión de la distancia en cada punto; la de Camera es la riqueza de densidad de píxeles y de información semántica. El juicio de Lang Xianpeng es que el LiDAR solo entrega una nube de puntos de alrededor de 128 líneas, mientras que la cámara entrega un mundo a color de varios millones de píxeles. La conducción autónoma no necesita solo distancias, sino también objetos, semántica, intenciones de movimiento y el contexto del entorno.
\end{knowledgebox}

\subsection{Qué resuelven BEV y Transformer}

Esta subsección explica BEV. BEV significa Bird's Eye View, vista de pájaro. El enfoque ingenuo consiste en que cada cámara reconozca primero los objetivos por su cuenta y luego se fusionen los resultados a posteriori; el problema es que las distintas cámaras pueden omitir detecciones, superponerse o tener escalas inconsistentes, y la fusión tardía es propensa a errores. La idea de BEV es extraer primero las características de las múltiples cámaras, unificar la información en una sola representación espacial y luego reconocer de una vez personas, vehículos, carriles y obstáculos. Aquí Transformer es una arquitectura de modelado de propósito general; lo importante es que está al servicio de la comprensión espacial unificada.

\begin{importantbox}{Lo esencial de BEV no es que la vista de pájaro se vea bien, sino que evita los errores de fusión tardía}
La fusión tardía es «cada cámara juzga primero de forma independiente y luego se ensamblan los resultados»; BEV se acerca más a «primero se reúnen las características de todas las cámaras y luego se juzga en un espacio unificado». Esto reduce los conflictos entre resultados de distintas fuentes y mejora la consistencia de los objetos.
\end{importantbox}

\subsection{Fusión temprana y fusión tardía: por qué la idea importa más que el operador}

Esta subsección amplía la explicación de Lang Xianpeng sobre BEV. En los primeros esquemas de varias cámaras, el enfoque ingenuo habitual era que cada cámara ejecutara su propia percepción hacia adelante: la cámara frontal reconocía los carriles y vehículos de enfrente, las cámaras laterales los objetivos a los costados y la cámara trasera los objetivos de atrás; después, el sistema de ingeniería fusionaba esas detecciones en un mismo sistema de coordenadas. El problema es que el resultado de cada cámara puede estar equivocado, incompleto, duplicado o tener una escala inconsistente, y el posprocesamiento debe decidir en qué cámara confiar, si dos mitades de vehículo son el mismo vehículo y cómo empalmar los bordes de las oclusiones. Esto consume una gran cantidad de esfuerzo en «eliminar los errores de la fusión tardía».

La elevación de dimensión de BEV consiste en adelantar la fusión al nivel de las características: primero se extraen características de las imágenes de las múltiples cámaras y luego se infieren objetos, carriles y zonas transitables en un espacio unificado. No se trata simplemente de «unir las imágenes en una imagen grande», sino de que el modelo use al mismo tiempo, en un mismo espacio, la información de varios puntos de vista. El docente enfatiza que Transformer es solo la estructura de red que después se puso al servicio de esta idea; lo verdaderamente decisivo fue haber identificado antes el problema de fondo de la industria de la conducción autónoma: la percepción desde múltiples puntos de vista debe modelarse de forma unificada, no juzgarse localmente en cada vista para luego ensamblarse.

\begin{knowledgebox}{Nota de clase: la lesson de BEV puede trasladarse a muchos sistemas de IA}
Cuando un sistema tiene varias fuentes de información, la fusión tardía suele traer problemas de consistencia. La solución de fondo es hacer una fusión temprana en un espacio de representación adecuado, para que el modelo aproveche todo el contexto de una sola vez. Los modelos multimodales, el estado de las herramientas de un Agent y los VLA de robótica siguen una lógica similar.
\end{knowledgebox}

\begin{figure}[H]
\centering
\includegraphics[width=0.92\textwidth]{av-terms-map.png}
\caption{Mapa de términos clave de la conducción autónoma: la relación entre BEV, Transformer, E2E, VLM, World Model y RL. Diagrama conceptual propio, elaborado a partir de la conversación de 00:20:06--01:26:40.}
\end{figure}

\begin{knowledgebox}{Lectura de la figura: los términos no forman una lista de elementos paralelos}
BEV y Transformer resuelven principalmente la representación y el modelado espaciales; E2E deja al modelo el mapeo de la entrada a la trayectoria; VLM añade comprensión semántica; World Model predice estados futuros; RL permite que el sistema se optimice a partir de la retroalimentación. Al leer esta figura, hay que verla como una cadena que va de la percepción a la decisión y luego al aprendizaje en lazo cerrado.
\end{knowledgebox}

\subsection{Por qué desarrollar chips propios}

Esta subsección trata los chips. El BEV con varias cámaras y el cómputo de modelos grandes requieren mucha capacidad de cómputo en el vehículo. En sus inicios, Tesla pasó de los chips de NVIDIA a su propio chip FSD con esta lógica: solo diseñando en conjunto el algoritmo, los sensores y el chip se obtiene capacidad de cómputo efectiva a un costo relativamente bajo. En la entrevista se menciona que el costo de sensores más chip de Tesla ronda el orden de 1,000 dólares; lo importante no es la cifra exacta, sino mostrar que una ruta de producción en serie debe diseñar de manera conjunta la capacidad de cómputo, el costo y la adaptación del algoritmo.

\begin{knowledgebox}{Términos clave: TOPS, ASIC, capacidad de cómputo efectiva}
\noindent\begin{tabularx}{\textwidth}{p{0.18\textwidth}p{0.36\textwidth}X}
\toprule
Término & Significado & Papel en este episodio \\
\midrule
TOPS & Billones de operaciones por segundo; indicador habitual de la capacidad de cómputo de los chips de IA para vehículos & Un valor alto no garantiza que el modelo real funcione bien. \\
ASIC & Circuito integrado de aplicación específica, hecho a la medida de un algoritmo o carga de trabajo concretos & Tesla puede adaptar mutuamente su algoritmo y su chip propios. \\
Chip de propósito general & Procesador que admite varias empresas y diversos algoritmos & Es flexible, pero no necesariamente óptimo para un algoritmo dado. \\
Capacidad de cómputo efectiva & Capacidad de cómputo que el modelo aprovecha realmente & Se acerca más al resultado de ingeniería que los TOPS nominales. \\
\bottomrule
\end{tabularx}
\end{knowledgebox}

\subsection{Resumen de la sección}

La esencia de la ruta de Tesla es llevar la conducción autónoma de la certeza basada en sensores costosos y mapas hacia la visión, la representación espacial unificada, los datos de la flota, los chips propios y el entrenamiento a gran escala. Fue la primera vez que la conducción autónoma se orientó explícitamente hacia un «sistema basado en la capacidad del modelo».
\section{De extremo a extremo: del desarrollo de funciones al entrenamiento de capacidades}

La sección anterior explicó cómo BEV unifica el espacio; esta sección examina un giro mayor: el enfoque de extremo a extremo. Lang Xianpeng (郎咸朋) dijo: «antes hacíamos mal la conducción autónoma». Se refería a desarrollar la conducción autónoma como si fuera una función: enumerar escenarios, escribir reglas, delimitar el ODD y corregir los bugs uno por uno con parches. Pero los escenarios de conducción reales son prácticamente imposibles de enumerar por completo, y las variables tampoco son independientes entre sí: corregir un escenario puede romper otro. El sentido del enfoque de extremo a extremo es convertir la conducción de «reglas escritas a mano» en «capacidades que el modelo aprende de los datos».

\begin{figure}[H]
\centering
\includegraphics[width=0.96\textwidth]{e2e-stack.png}
\caption{Pila de conducción autónoma de extremo a extremo: de la entrada de sensores a la salida de trayectoria, las capas intermedias se van comprimiendo en un sistema que aprende. Diagrama conceptual propio, elaborado a partir de la conversación entre 00:28:46 y 00:41:14.}
\end{figure}

\begin{knowledgebox}{Lectura de la figura: de extremo a extremo no significa sin estructura, sino que el modelo asume más correspondencias}
De Sensors, BEV, Planning y Trajectory hasta Control, en el sistema tradicional cada capa tiene una gran cantidad de reglas escritas a mano. El enfoque de extremo a extremo no elimina por completo la ingeniería, sino que entrega las correspondencias clave al modelo y a los datos, para que el sistema aprenda la capacidad de ir de la entrada a la trayectoria.
\end{knowledgebox}

\subsection{Por qué no es viable enumerar todos los escenarios}

Esta subsección desglosa el fracaso de los sistemas basados en reglas. Las combinaciones de variables como el clima, el tráfico, la estructura de la vía, el estado del pavimento, los vehículos cercanos, los peatones, la velocidad del propio vehículo, las obras, los baches y los objetos extraños son enormes. Aun si se definieran a mano decenas de millones de escenarios, surgirían dos problemas: los escenarios no son ortogonales y las reglas se afectan entre sí; y siempre aparecerán eventos de cola larga que una persona evita con sentido común, pero que un sistema de reglas quizá no sepa cómo manejar.

\begin{warningbox}{El dilema de la cola larga en los sistemas de reglas}
La conducción autónoma no es una lógica de botones. Una regla de vuelta a la derecha se comporta de forma completamente distinta con sol, con lluvia, con tráfico congestionado, con peatones, con bicicletas, con obras o con baches en el pavimento. Cuantas más reglas hay, más graves son las interferencias entre ellas y mayor es el costo de mantenimiento del sistema.
\end{warningbox}

\subsection{Un sistema de capacidades: como enseñar a un niño, no aprender por él}

Esta subsección conserva una metáfora oral muy importante de Lang Xianpeng. Con el enfoque de extremo a extremo, el ingeniero es como un observador y un maestro: proporciona datos de alta calidad, recursos de entrenamiento, preguntas de evaluación y retroalimentación, pero no puede aprender en lugar del modelo. El ingeniero no necesariamente puede explicar del todo en qué se convierte el modelo ni cómo se forman sus estrategias internas; pero puede evaluar si progresa mediante case, pruebas de regresión, simulación, pruebas en carretera y retroalimentación de los usuarios.

\begin{importantbox}{Nota de clase: el modelo y los datos son dos caras de la misma moneda}
Construir un sistema de capacidades con un modelo implica necesariamente depender de los datos; solo las empresas que poseen datos de alta calidad, verticales y con retroalimentación real tienen la oportunidad de mejorar de forma continua la conducción autónoma. Los parámetros del modelo, las técnicas algorítmicas y la experiencia del equipo importan, pero sin datos de conducción de alta calidad el techo será muy bajo.
\end{importantbox}

\subsection{Qué hacen los ingenieros después del enfoque de extremo a extremo}

Esta subsección aclara un punto que se presta a malentendidos. El enfoque de extremo a extremo no significa que los ingenieros se retiren, ni que la responsabilidad de seguridad se entregue a una caja negra. Al contrario, el trabajo del ingeniero pasa de «escribir reglas para cada escenario» a «definir los datos, la evaluación, la regresión y los límites de seguridad». Deben encontrar los case de cola larga y juzgar qué escenarios revelan los puntos ciegos del modelo; deben diseñar el etiquetado y los datos de entrenamiento para que el modelo vea muestras de alto valor; deben construir la simulación y la evaluación en lazo cerrado para evitar que un modelo nuevo corrija un escenario y rompa otro; y deben controlar el ritmo de liberación, para asegurar que los usuarios de vehículos de producción en serie no sean sujetos de experimentación sin protección.

\begin{knowledgebox}{Experiencia práctica: el enfoque de extremo a extremo traslada la dificultad de ingeniería a los datos y la evaluación}
En la era de las reglas, la dificultad de ingeniería estaba en las ramas de código y en la enumeración de escenarios; en la era de los modelos, está en la cobertura de datos, el peso de las muestras, la estabilidad del entrenamiento, las métricas fuera de línea, la experiencia en línea y la regresión de seguridad. El trabajo de ingeniería no desaparece: cambia de lugar.
\end{knowledgebox}

\noindent\begin{tabularx}{\textwidth}{p{0.22\textwidth}p{0.34\textwidth}X}
\toprule
Rol de ingeniería & Trabajo principal en la era de las reglas & Trabajo principal en la era de extremo a extremo \\
\midrule
Ingeniero de algoritmos & Escribir detección, planificación, reglas y máquinas de estados & Diseñar la estructura del modelo, los objetivos de entrenamiento, el muestreo de datos y la evaluación. \\
Ingeniero de datos & Recopilar unos pocos case para validar reglas & Construir las cadenas de procesamiento para el retorno de datos de la flota, la limpieza, la anonimización, el etiquetado y la minería. \\
Ingeniero de pruebas & Ejecutar escenarios fijos y regresiones manuales & Mantener la biblioteca de case de cola larga, la simulación, las métricas fuera de línea y el despliegue gradual en línea. \\
Producto/seguridad & Definir los límites de la función y los avisos & Definir la responsabilidad de la toma de control, el ODD, la degradación ante riesgos y la educación del usuario. \\
\bottomrule
\end{tabularx}

\subsection{De extremo a extremo, RL, VLM y modelos del mundo}

Esta subsección se conecta con la ruta posterior. En la entrevista se discuten juntos el enfoque de extremo a extremo, los VLM, los modelos del mundo y el RL. El enfoque de extremo a extremo hace primero que el sistema aprenda el comportamiento de conducción desde la percepción hasta la trayectoria; los VLM incorporan comprensión semántica e interpretación de escenarios complejos; los modelos del mundo predicen el efecto de las acciones sobre el estado futuro del tráfico; y el RL aporta un marco de optimización en lazo cerrado. La conducción autónoma se parece cada vez más a un Agent en el mundo físico: observar, predecir, actuar, recibir retroalimentación y volver a aprender.

\begin{figure}[H]
\centering
\includegraphics[width=0.96\textwidth]{e2e-rl-worldmodel.png}
\caption{De extremo a extremo + VLM + modelo del mundo: de la imitación de trayectorias al aprendizaje por refuerzo y la predicción de estados futuros. Diagrama conceptual propio, elaborado a partir de la conversación entre 01:15:15 y 01:26:40.}
\end{figure}

\begin{knowledgebox}{Lectura de la figura: lo que aporta el modelo del mundo es la «predicción de consecuencias»}
Una política de extremo a extremo puede producir una trayectoria, pero el tráfico complejo exige predecir: si cambio de carril, reduzco la velocidad o cedo el paso, cómo cambiarán los vehículos y las personas alrededor. El modelo del mundo incorpora al sistema las consecuencias de las acciones, y el RL permite que el sistema optimice su política a partir de la retroalimentación.
\end{knowledgebox}

\subsection{Resumen de la sección}

La esencia del enfoque de extremo a extremo no es una palabra de moda, sino el paso de la conducción autónoma del desarrollo de funciones de software al entrenamiento de capacidades. Convierte al ingeniero, de alguien que escribe reglas a mano, en diseñador de los datos, el entrenamiento, la evaluación y los límites de seguridad.
\section{L2, L3 y L4: los límites de responsabilidad importan más que el número de nivel}

La sección anterior explicó el entrenamiento de capacidades; esta sección examina los niveles de conducción autónoma. Lang Xianpeng (郎咸朋) enfatiza que «L3 no es una extensión de L2, sino el precursor de L4». Lo decisivo de este juicio no está en el número, sino en la responsabilidad: L2 es asistencia a la conducción, y el conductor humano es responsable; en L3, bajo condiciones específicas, el sistema empieza a asumir la responsabilidad y avisa a la persona con anticipación cuando necesita que tome el control; L4, en cambio, es aquel en que el sistema es plenamente responsable dentro de una zona delimitada o de escenarios delimitados.

\begin{figure}[H]
\centering
\includegraphics[width=0.92\textwidth]{l3-l4-ladder.png}
\caption{L3 no es una extensión de L2: L3 se parece más a un precursor de L4 que a una mejora lineal de la asistencia a la conducción. Diagrama conceptual propio, elaborado a partir de la conversación entre 00:50:50 y 01:15:15.}
\end{figure}

\begin{knowledgebox}{Lectura de la figura: detrás del cambio de nivel hay un cambio de responsabilidad}
A la izquierda, en L2, la persona es responsable y el sistema asiste; a la derecha, L3/L4 empiezan a exigir que el sistema asuma la responsabilidad y que se definan con claridad el ODD y los límites de la toma de control. Al leer los niveles, no basta con fijarse en cuántas funciones hay: hay que considerar la responsabilidad en caso de accidente, el mecanismo de toma de control, la redundancia de seguridad y los límites de operación.
\end{knowledgebox}

\subsection{Por qué L2 no puede convertirse de forma natural en L3}

Esta sección explica un error de concepto. Muchas personas imaginan L2, L3 y L4 como niveles continuos: cada vez más funciones y una experiencia cada vez mejor. Pero, desde el punto de vista de la responsabilidad, el paso de L2 a L3 es un cambio cualitativo. L2 puede exigir que la persona esté atenta en todo momento y tome el control; L3, en cambio, debe garantizar seguridad suficiente durante el periodo en que el sistema es responsable y, al salir del ODD o de los límites de su capacidad, ofrecer una ventana de toma de control que la persona pueda ejecutar. Esto involucra la regulación, la definición del producto, la estrategia de seguridad y la educación de los usuarios.

\begin{warningbox}{Cuanto mejor es L2, más fácil es que el usuario lo use mal}
Si la experiencia de L2 se acerca a la conducción autónoma, pero la responsabilidad legal y de seguridad sigue recayendo en la persona, el usuario puede confiar en exceso en el sistema. Este es el punto más peligroso del límite entre L2 y L3: la experiencia parece de conducción autónoma, pero la responsabilidad sigue siendo de la persona.
\end{warningbox}

\subsection{ODD y toma de control: cómo se traslada el límite de responsabilidad al producto}

Esta sección lleva la cuestión de los niveles al diseño del producto. L3/L4 no consiste en anunciar un nivel más en una presentación, sino en que el sistema sepa dónde puede asumir la responsabilidad, cuándo debe salir, cómo avisar a la persona para que tome el control y cómo pasar a un estado de riesgo mínimo si la toma de control falla. Por ejemplo, las autopistas, las vías rápidas urbanas, el seguimiento de vehículos en congestión, el estacionamiento, ciertas condiciones climáticas y ciertos rangos de velocidad pueden constituir ODD distintos. En cuanto el sistema se compromete a ser responsable dentro de un ODD, necesita el diseño correspondiente de redundancia de percepción, diagnóstico de fallas, monitoreo de estado y responsabilidad regulatoria.

\begin{importantbox}{Tres preguntas sobre el límite de responsabilidad}
Primera: ¿cuándo puede el sistema asumir la responsabilidad de la conducción? Segunda: ¿cuándo debe el sistema devolver la responsabilidad a la persona o pasar a un estado degradado? Tercera: si la persona no toma el control a tiempo, ¿cómo detiene el sistema el vehículo de forma segura o mantiene el riesgo al mínimo? Mientras estas tres preguntas no tengan una respuesta clara, L3 es solo un término de mercadotecnia.
\end{importantbox}

\subsection{Por qué L3 es el precursor de L4}

Esta sección explica el sentido de «precursor». L3 exige que el sistema sea realmente responsable bajo ciertas condiciones, y eso obliga a la empresa a construir las capacidades que necesita L4: definición del ODD, redundancia de seguridad, estrategia de toma de control, minimización del riesgo, monitoreo del sistema, responsabilidad en caso de accidente y un sistema de validación. Por eso no es solo una mejora del paquete de funciones de L2, sino el campo de entrenamiento para entrar en la etapa en que «el sistema es responsable».

\noindent\begin{tabularx}{\textwidth}{p{0.16\textwidth}p{0.36\textwidth}X}
\toprule
Nivel & Responsabilidad central & Implicación de ingeniería \\
\midrule
L2 & La persona es responsable; el sistema asiste & La experiencia funcional puede ser muy buena, pero la persona debe supervisar. \\
L3 & El sistema es responsable bajo condiciones específicas; requiere un mecanismo de toma de control & Se vuelven decisivos los límites de responsabilidad, la estrategia de salida y la certificación regulatoria. \\
L4 & El sistema es plenamente responsable en zonas o escenarios delimitados & Requiere un ciclo cerrado de seguridad completo y la gestión del dominio de operación. \\
L5 & Conducción autónoma sin restricciones & Por ahora se parece más a un concepto de largo plazo. \\
\bottomrule
\end{tabularx}

\subsection{Resumen de la sección}

Los niveles de conducción autónoma no son números de mercadotecnia, sino una estructura de responsabilidad. El valor principal de L3 es que obliga a la empresa a enfrentar las cuestiones de seguridad, regulación, toma de control y responsabilidad del sistema que requiere L4.
\section{La experiencia de Li Auto: datos, presión organizacional y el proyecto Weicheng}

Las secciones anteriores trataron las rutas técnicas; esta sección vuelve a la práctica de Li Auto (理想汽车). Cuando Lang Xianpeng (郎咸朋) se incorporó a Li Auto en 2018, ya había discutido con Li Xiang (李想) que lo más importante en la conducción autónoma son los datos: el talento se puede reclutar y la capacidad de cómputo se puede comprar, pero los datos reales de la flota, los casos de cola larga y los escenarios de los usuarios no se pueden comprar fuera. La ventaja de un fabricante de automóviles de producción en serie está en tener flota, usuarios y retroalimentación de caminos reales; la dificultad es que debe iterar rápido bajo presiones de seguridad, costo, experiencia y organización.

\begin{figure}[H]
\centering
\includegraphics[width=0.96\textwidth]{ideal-av-data-loop.png}
\caption{Ciclo cerrado de datos de la conducción inteligente de Li Auto: la flota de producción en serie, los casos de cola larga, el entrenamiento y la retroalimentación del despliegue forman un volante de inercia. Diagrama conceptual propio, elaborado a partir de la conversación de 00:41:14--00:50:50 y 01:26:40--01:35:51.}
\end{figure}

\begin{knowledgebox}{Lectura de la figura: la ventaja de datos del fabricante proviene de un ciclo cerrado real}
De Fleet, Cases, Label/Train y Deploy a Feedback. La flota de producción en serie no es un almacén de datos, sino un ciclo cerrado que descubre continuamente problemas de cola larga, entrena modelos, los despliega por OTA y vuelve a recoger retroalimentación. El techo de la conducción autónoma lo determinan conjuntamente la calidad de los datos y la capacidad de cómputo para entrenamiento.
\end{knowledgebox}

\subsection{Privacidad de datos y datos del exterior del vehículo}

Esta sección agrega un límite importante. En la entrevista, Lang Xianpeng enfatizó que Li Auto recoge datos del exterior del vehículo y no recoge información biométrica de los usuarios, como rostros o voces dentro del vehículo; además, al transmitir los datos del exterior se procesan los rostros y las placas. Esto no es un detalle de relaciones públicas, sino la base para que el ciclo cerrado de datos de la conducción autónoma pueda operar a largo plazo: sin cumplimiento normativo y sin la confianza de los usuarios, el volante de inercia de datos reales no puede sostenerse.

\begin{importantbox}{Experiencia práctica: el ciclo cerrado de datos requiere confianza como condición previa}
La conducción autónoma necesita una gran cantidad de datos reales de caminos, pero los datos reales no son un recurso sin límites. La captura en el exterior del vehículo, la anonimización, el control de permisos, la limitación de usos y el almacenamiento seguro determinan si los datos de la flota pueden convertirse en un activo de largo plazo.
\end{importantbox}

\subsection{De Baidu a Li Auto: cómo la formación en una ruta cambia la organización}

Esta sección complementa la trayectoria personal de Lang Xianpeng. En 2013 participó en Street View y en mapas de alta precisión dentro del equipo de Baidu Maps, y después en la colaboración de conducción autónoma con BMW; en esa etapa se formaron capacidades en mapas, topografía, percepción y sistemas de vehículos de ingeniería. Al incorporarse a Li Auto en 2018 se enfrentó a un fabricante de producción en serie: los vehículos se venden a usuarios reales, las funciones deben llegar por OTA a una gran cantidad de vehículos y los problemas regresan de los caminos reales en forma de casos de cola larga. La diferencia entre estas dos etapas explica por qué insiste una y otra vez en los datos y en el ciclo cerrado de la producción en serie.

\begin{knowledgebox}{Nota de clase: el vehículo de investigación, el Robotaxi y el vehículo de producción en serie son tres organizaciones distintas}
El vehículo de investigación busca la validación técnica; el Robotaxi busca poder operar dentro de una zona delimitada; el vehículo de producción en serie busca costo, estabilidad, experiencia de usuario y mantenimiento durante todo su ciclo de vida. Las rutas técnicas son parecidas, pero los objetivos organizacionales son completamente distintos.
\end{knowledgebox}

\noindent\begin{tabularx}{\textwidth}{p{0.18\textwidth}p{0.36\textwidth}X}
\toprule
Etapa & Activos clave & Formación para la conducción autónoma \\
\midrule
Baidu Maps/Street View & Vehículos de captura, topografía, anonimización, mapas de alta precisión & Comprender los datos de caminos y la infraestructura de mapas. \\
Baidu ADU/colaboración con BMW & Vehículos de ingeniería, pruebas L3/L4, mapas de alta precisión & Conocer la ruta determinista temprana y las limitaciones de los vehículos de ingeniería. \\
Vehículos de producción en serie de Li Auto & Flota de usuarios, OTA, datos del exterior del vehículo, recursos organizacionales & Construir un ciclo cerrado de datos reales y la conciencia de responsabilidad de la producción en serie. \\
\bottomrule
\end{tabularx}

\subsection{El proyecto Weicheng y la presión organizacional}

Esta sección aborda la historia organizacional de la segunda mitad de la entrevista. El proyecto «Weicheng» (卫城) se relata como una experiencia de alta presión: la negociación con proveedores, la elección de ruta técnica, los plazos de producción en serie, la presión de Li Xiang y la dignidad del equipo se entrelazan. Cuando Lang Xianpeng habla de «morir de pie», lo que hay detrás es un equipo de tecnología de frontera que, en un momento decisivo, no quiso asegurar la seguridad de corto plazo arrodillándose ante proveedores externos, sino que eligió asumir el riesgo organizacional de cambiar de ruta.

\begin{figure}[H]
\centering
\includegraphics[width=0.86\textwidth]{weicheng-project.png}
\caption{Campo de presiones del proyecto Weicheng: la ruta técnica, los proveedores, la dignidad organizacional y los plazos de producción en serie presionan al mismo tiempo. Diagrama conceptual propio, elaborado a partir de la conversación de 01:32:23--01:35:51.}
\end{figure}

\begin{knowledgebox}{Lectura de la figura: la elección de ruta técnica se convierte en presión organizacional}
En el centro de la figura está Project Weicheng, rodeado por tecnología, proveedores, organización, producción en serie, liderazgo y equipo. Al leer esta figura hay que ver que la conducción autónoma no es una competencia de modelos de laboratorio: la elección de ruta afecta la entrega de la producción en serie, la relación con la cadena de suministro, la moral del equipo y la estrategia de la empresa.
\end{knowledgebox}

\subsection{Liderazgo de alta presión y ejecución de ingeniería}

Esta sección trata con cautela la discusión relacionada con Li Xiang. La entrevista menciona que Li Xiang se enojaba con el equipo de conducción inteligente y que no tiene muy buen carácter, y también menciona el objetivo detrás de esa presión: que el equipo completara la transición técnica y la entrega de la producción en serie dentro de una ventana de tiempo crítica. Al organizar este material como curso no se evalúa el carácter de la persona, sino el mecanismo organizacional: si la alta presión trae objetivos más claros, una asignación de recursos más rápida y una ejecución más sólida, y si al mismo tiempo evita el desgaste excesivo, la cultura del miedo y las concesiones en seguridad.

\begin{warningbox}{La alta presión no es correcta por sí misma}
Las organizaciones de tecnología de frontera necesitan objetivos firmes y una ejecución sólida, pero la conducción autónoma implica responsabilidad sobre la seguridad, y la alta presión no puede sustituir la validación rigurosa. Una buena presión organizacional debe impulsar que los problemas salgan a la luz, que los recursos se concentren y que las decisiones se tomen rápido, no obligar al equipo a saltarse los límites de seguridad.
\end{warningbox}

\subsection{Negociación con proveedores y capacidad de desarrollo propio}

Esta sección sitúa las expresiones «arrodillarse a suplicarle» y «morir de pie» en su contexto industrial. Los fabricantes de automóviles con conducción autónoma suelen depender de proveedores de chips, sensores, algoritmos, mapas, cadenas de herramientas e integración. Depender de proveedores puede acelerar la implementación temprana, pero si las capacidades centrales se subcontratan durante mucho tiempo, en los momentos decisivos el fabricante pierde la iniciativa sobre su ruta. La emoción con la que Lang Xianpeng habla del proyecto Weicheng no es solo cuestión de carácter personal, sino de un equipo que elige entre «seguir dependiendo de soluciones externas» y «arriesgarse a construir capacidad de desarrollo propio».

\begin{importantbox}{Las capacidades centrales no pueden depender solo de compras}
La barrera de largo plazo de la conducción autónoma está en el ciclo cerrado de datos, modelos, evaluación, despliegue en el vehículo y responsabilidad sobre la seguridad. Los proveedores externos pueden aportar herramientas y componentes, pero si el fabricante no domina el criterio central y la capacidad de iteración, le será muy difícil cambiar de rumbo con rapidez en la era de extremo a extremo.
\end{importantbox}

\begin{figure}[H]
\centering
\includegraphics[width=0.96\textwidth]{av-org-leadership.png}
\caption{Ciclo cerrado organizacional de la conducción autónoma: el criterio técnico, la presión de la dirección, la ejecución del equipo y la responsabilidad sobre la seguridad deben cumplirse al mismo tiempo. Diagrama conceptual propio, elaborado a partir de la conversación de 01:26:40--02:00:20.}
\end{figure}

\begin{knowledgebox}{Lectura de la figura: el ciclo cerrado organizacional es tan importante como el técnico}
El equipo de conducción autónoma necesita juzgar la ruta técnica, soportar la presión de la alta dirección, completar la ejecución de ingeniería y después construir confianza mediante la evaluación de seguridad y la experiencia de usuario. Si cualquier eslabón se rompe, a la ruta técnica le será muy difícil llegar al vehículo de producción en serie.
\end{knowledgebox}

\subsection{Resumen de la sección}

La experiencia de Li Auto muestra que la competencia entre empresas de conducción autónoma no se limita a modelos y algoritmos, sino que también abarca el cumplimiento normativo de los datos, el ciclo cerrado de la flota, la negociación en la cadena de suministro, la resiliencia organizacional y la validación de seguridad. La conducción inteligente de producción en serie se forma con el entrenamiento conjunto de un sistema técnico y un sistema organizacional.
\section{Síntesis y ampliación}

Esta sección condensa todo el episodio en una hoja de ruta. La evolución de la conducción autónoma a lo largo de diez años puede verse como tres «aumentos de dimensión»: el primero va de las reglas y los mapas de alta precisión a la representación espacial unificada con BEV/Transformer; el segundo, del desarrollo modular de funciones al entrenamiento de capacidades de extremo a extremo; el tercero, de la imitación de las trayectorias de conducción humanas al Agent del mundo físico basado en VLM, modelos del mundo y RL.

\begin{importantbox}{Seis conclusiones centrales}
Primera: la ruta de mapas de alta precisión y LiDAR resuelve la certidumbre a corto plazo, pero difícilmente escala al nivel de ciudades enteras. Segunda: la Camera tiene alta densidad de información y bajo costo, y es la base de la ruta de visión para producción en serie. Tercera: lo esencial de BEV es unificar la representación espacial de varias cámaras y reducir los errores de fusión tardía. Cuarta: el enfoque de extremo a extremo representa un cambio de paradigma, del desarrollo de funciones al entrenamiento de capacidades. Quinta: la esencia de L3/L4 es el límite de responsabilidad, no una lista de funciones. Sexta: la barrera de largo plazo de la conducción autónoma en producción en serie proviene de los datos reales, la capacidad de cómputo, la ejecución organizacional y el ciclo cerrado de validación de seguridad.
\end{importantbox}

\subsection{Tabla de consulta rápida de términos clave}

\noindent\begin{tabularx}{\textwidth}{p{0.22\textwidth}p{0.34\textwidth}X}
\toprule
Término & Significado en este episodio & Relación con el curso \\
\midrule
Mapa de alta precisión & Registro previo de la geometría y la información semántica de las vías & Dependencia central de la primera ruta basada en certidumbre. \\
LiDAR & Sensor activo de medición de distancia & Da geometría directa, pero tiene restricciones de costo y de densidad de información. \\
Camera & Sensor de visión pasivo & Alta densidad de información; adecuado para producción en serie de bajo costo y datos a gran escala. \\
BEV & Bird's Eye View, representación espacial a vista de pájaro & Paso clave para el modelado unificado de varias cámaras. \\
Transformer & Arquitectura general de modelado de secuencias/atención & Se trasladó a la comprensión espacial en la conducción autónoma. \\
De extremo a extremo & Mapeo modelado desde la entrada hasta la trayectoria/acción & La conducción autónoma pasa de la ingeniería de funciones al entrenamiento de capacidades. \\
VLM & Vision-Language Model, modelo de visión y lenguaje & Introduce comprensión semántica e interpretación de escenarios complejos. \\
World Model & Modelo que predice las consecuencias de las acciones y el estado futuro del mundo & Sustenta la planificación, la simulación y el ciclo cerrado de RL. \\
RL & Reinforcement Learning, aprendizaje por refuerzo & Optimiza la política con retroalimentación; es una de las direcciones de la conducción autónoma en ciclo cerrado. \\
\bottomrule
\end{tabularx}

\subsection{Preguntas para seguir observando}

\begin{enumerate}
\item ¿Podrá la ruta de extremo a extremo superar de forma estable a los sistemas basados en reglas en más ciudades, condiciones climáticas, estructuras viales y escenarios de cola larga?
\item ¿Cómo se reequilibrarán las rutas Camera-only, LiDAR-heavy y de fusión multisensor entre costo, seguridad y regulación?
\item ¿Se convertirá el modelo del mundo en el núcleo de la siguiente etapa de la conducción autónoma, o solo servirá como apoyo para entrenamiento/simulación?
\item Cuando L3 llegue a la producción en serie en China, ¿cómo deben diseñarse el límite de responsabilidad, el mecanismo de toma de control y la educación del usuario?
\item ¿Puede el ciclo cerrado de datos de flotilla de los fabricantes de automóviles convertirse en una barrera de largo plazo difícil de replicar para las empresas de modelos?
\item ¿Cómo equilibrar la alta presión organizacional y la validación de seguridad para no sacrificar la confiabilidad en aras de la velocidad de lanzamiento?
\end{enumerate}

\subsection{Lecturas adicionales}

\begin{itemize}
\item Si le interesan los modelos de negocio de la conducción autónoma y las diferencias entre Waymo y Momenta, puede contrastarlo con la entrevista del EP132 a Gao Jiyang (高继扬), de Xinghaitu (星海图).
\item Si le interesan la Physical AI y la fábrica de modelos grandes a bordo del vehículo, puede contrastarlo con la entrevista del EP120 a Liu Xianming (刘先明), de Xiaopeng (小鹏).
\item Si le interesan VLA, los modelos del mundo y los modelos fundacionales para robótica, puede contrastarlo con la versión de pantalla compartida sobre VLA `eiQFomOuCJs` y con la entrevista sobre multimodalidad del EP102 a Zhang Xiangyu (张祥雨).
\item Si le interesan los costos de hardware y la industria del LiDAR, puede contrastarlo con la entrevista del EP111 a Hesai (禾赛).
\end{itemize}

\end{document}
